\section{The Fear That Grows with Love}

Fear of punishment can begin to loosen the hold of sin, while fear of offending God deepens as love becomes stronger. St.~Thomas distinguishes these concerns in his treatment of the gift of fear. A person may first refrain because he dreads the penalty and then become more attentive to the offense against the One he loves. The concerns can remain together while charity grows. His account explains how an imperfect beginning can prepare the will for a fuller response to God, and why increasing love need not make the Christian careless about salvation.

Thomas begins by asking in what sense God is to be feared. God is goodness itself; we do not fear him as though he were an evil. We rightly fear his just punishment, and we fear being separated from him through our fault. His justice gives reason for fear, while his mercy gives reason for hope (\emph{Summa theologiae} II--II, q.~19, a.~1). The two movements can therefore belong to the same soul before the same Lord.

The fear called servile regards punishment. The fear called filial regards offending God, as a child fears to offend a beloved father (a.~2). The names identify different objects of concern. In the first, the evil that arrests attention is what sin will bring upon the sinner. In the second, the heart has become attentive to the wrong of turning against the One it loves. Filial fear thus includes a personal relation to God, received through charity, and a desire to remain faithfully within that relation.

Thomas teaches that fear of punishment is good in itself, while distinguishing it from the servility which resists love. A person may refrain from a crime because he dreads its penalty, yet still cherish the injustice and wish he could commit it safely. The defect lies in the refusal to love justice, not in recognizing that punishment is dreadful (a.~4). A meditation on hell should welcome the warning's power to restrain sin and ask that the heart itself be changed. Mere outward restraint leaves work for grace to accomplish within.

Nor must a person prove that he loves God by becoming indifferent to his own salvation. Thomas explains that concern for one's own good can be ordered within charity. Fear of punishment can remain alongside love when the person's own comfort has not become his supreme end. Indeed, charity intensely rejects the loss of God, which is itself the chief punishment (a.~6). We were made for communion with him. To desire that communion and dread its loss is consistent with loving him above all things.

This distinction brings clarity to a conscience troubled by mixed motives. A penitent can be sorry that he offended God and also afraid of the punishment he deserved. The presence of the latter concern does not by itself make the former unreal. He should turn toward God with the contrition he has, seek the sacrament, and ask for a more generous love. An endless effort to separate every movement of the heart can become another way of delaying obedience.

Thomas calls fear a beginning of wisdom in more than one way. Fear of punishment can first restrain the sin that obstructs a wise life. Filial fear belongs to wisdom already beginning its work: the person reveres God and submits to being ruled by him (a.~7). The difference becomes visible in ordinary conduct. Someone may initially stop a dishonest act because judgment has become frightening. As love matures, he also wants the truth to govern his dealings because they are offered to God.

The gift of fear belongs to this filial disposition. Thomas describes a soul becoming responsive to the Holy Spirit through reverence and through refusal to separate itself from God (a.~9). This has a quiet firmness. The person becomes more willing to be corrected, less determined to reserve one protected area against God's command, and more ready to acknowledge dependence. The Lord's nearness does not make him ordinary or place him at our disposal. Familiarity in prayer can deepen awe when the soul remembers whose friendship it has received.

Filial fear and hope strengthen one another. Hope relies upon God's help; filial fear guards against withdrawing from that help (a.~9, reply~1). The soul does not become more hopeful by pretending that its own refusal is harmless. It becomes more hopeful by trusting the giver of grace and more watchful because it does not want to resist him. This union corrects both the arrogance which presumes upon pardon and the discouragement which treats amendment as impossible.

Thomas teaches that filial fear grows as charity grows, even while the active fear of punishment diminishes. Greater love makes offense and separation more hateful; closer adherence to God gives greater confidence (a.~10). This growth may appear in a quicker willingness to accept correction. A fault once defended has become grievous because the person more clearly recognizes whom it offends. He also has greater reason to seek help promptly: his hope rests upon the God whose friendship he values. Thomas's account explains how reverence and confidence can increase together in the same person.

A person beginning with fear of punishment can therefore pray for charity without pretending that he has ceased to dread the punishment. He should use the motive which restrains him and ask to love God more fully. Thomas's distinction prevents him from either despising the help already received or treating outward restraint as the whole of conversion. In the next temptation, the good first required may be simply to refuse consent. As the heart is healed, that refusal can become more willing because the person increasingly desires the friendship which sin would destroy.

